\subsection{The large-head chart by two Hermite bounds}
\label{sec:fqe-large-head-dual}
Throughout this subsection $q=1587/1000$, $r=q-1$, $L=\ln2$ and
$\theta=(397/100,3/20,159/100)$. Put
\begin{gather*}
 f(a,z)=\Lambda_\theta\left(a,\frac{(1-a)(1+2z)}3,
              \frac{(1-a)(1-z)}3,\frac{(1-a)(1-z)}3\right),\\
 B=\frac{1+2z^2}3,\qquad C=\frac{1+6z^2+2z^3}9,\qquad
 K=\frac{(1+2z)^q+2(1-z)^q}{3^q}.
\end{gather*}
Here $1/2\le a\le4/5$, $0\le z\le1$, and the purity and third moment are
$p=a^2+B(1-a)^2$, $t=a^3+C(1-a)^3$.
The objective is $f<T=2507333/1500000$. The first lemma supplies the
shape condition $f_{aaa}\le0$. The second then bounds $f$ by a quadratic
in $a$, replaces $K(z)$ by one cubic minorant, and proves positivity of
that quadratic's gap by completing the square. Thus the two-dimensional
chart reduces to univariate polynomial signs.

\begin{lemma}[Largest-entry curvature; exact polynomial sign certificates]\label{lem:fqe-large-head-curvature}
On this whole chart, $f_{aaa}\le0$.
\end{lemma}
\begin{proof}
Write $h(a,K)=-\ln(a^q+K(1-a)^q)/(rL)$. The remaining part of $f$ is a polynomial in $a$
whose only fourth derivative comes from $-\gamma p^2$, so
$f_{aaaa}=h_{aaaa}-24\gamma(1+B)^2$.

\emph{The shape of $h_{aaaa}$ in $K$.}
Set $x=2a-1\in[0,3/5]$, $\rho=(1-a)/a=(1-x)/(1+x)\in[1/4,1]$, $w=K\rho^q$ and
\[
 y=\frac{a^q-K(1-a)^q}{a^q+K(1-a)^q}=\frac{1-w}{1+w}.
\]
Apart from the positive factor $16q/(rL(1-x^2)^4)$, $h_{aaaa}$ is
\begin{align*}
 J(y)={}&6q^3y^4-24q^2xy^3+(-8q^3+36qx^2+8q)y^2\\
 &+(24q^2x-24x^3-24x)y+2q^3-36qx^2-8q+6x^4+36x^2+6.
\end{align*}
On $0\le x\le3/5$,
\[
 J'(y)\big|_{y=x}=8x(q-1)\{3(q-1)^2x^2-2q^2+q+3\}\le0,\qquad
 J'''(y)=144q^2(qy-x)\ge0\quad(y\ge x).
\]
So on $[x,1]$ the derivative $J'$ is convex and starts nonpositive. Hence $J$ decreases and then
increases, and has no interior maximum on any subinterval:
\[
 J(y)\le\max\{J(y_\ell),J(y_h)\}\qquad(x\le y_\ell\le y\le y_h\le1).
\]

\emph{Envelopes for $y$.}
Write $K=\lambda_1^q+2\lambda_2^q$ with $\lambda=\bigl(\frac{1+2z}3,\frac{1-z}3,\frac{1-z}3\bigr)$.
Convexity of $s^q$ gives $K\ge3^{-r}>1/2$, since $3^r<2$.
Concavity of $s^r$ with weights $\lambda_i$ gives $K=\sum\lambda_i\lambda_i^r\le B^r\le1$.
If $B\le1/2$, then $K\le2^{-r}<2/3$, because $2^q>3$.
For $0<\rho\le1$ we have $\rho^2\le\rho^q\le\rho^{3/2}\le\rho(1+\rho)/2$ and $\rho^q\le\rho$.
Hence $w\ge\rho^2/2$, $w\le\rho(1+\rho)/2=(1-x)/(1+x)^2$ always, and $w\le2\rho/3$ when $B\le1/2$.
Since $y=(1-w)/(1+w)$ decreases in $w$, we get $y_\ell\le y\le y_h$ with
\[
 y_h=\frac{(1+x)^2-(1-x)^2/2}{(1+x)^2+(1-x)^2/2},\quad
 y_A=\frac{(1+x)-\frac23(1-x)}{(1+x)+\frac23(1-x)},\quad
 y_B=\frac{(1+x)^2-(1-x)}{(1+x)^2+(1-x)},
\]
and $y_\ell=y_A$ if $B\le1/2$, $y_\ell=y_B$ otherwise.
Moreover $y_A-x$ and $y_B-x$ have the numerators $\frac13(1-x^2)$ and $(1-x)(x+x^2)$ over the
positive denominators, so $x\le y_\ell$, and $y_h\le1$.

\emph{Three polynomial inequalities.}
Let $L_0=0.693147<L$. For $0\le x\le3/5$,
\begin{align*}
 J(x,y_A)&\le\tfrac{1696}{25}\,\frac{rL_0(1-x^2)^4}{16q},\qquad
 J(x,y_h)\le\tfrac{1696}{25}\,\frac{rL_0(1-x^2)^4}{16q},\\
 J(x,y_B)&\le\tfrac{4293}{50}\,\frac{rL_0(1-x^2)^4}{16q}.
\end{align*}
Clearing the denominators ($D_A^4$ with $D_A$ linear, $D_B^4$ and $D_h^4$ with quadratic denominators)
turns each into a polynomial inequality of degree $12$, $16$ and $16$ in $x$.
Each is certified by the Bernstein expansion on $[0,3/5]$, all of whose coefficients are positive (LHD-1).
Thus $h_{aaaa}<1696/25=24\gamma(16/9)$ if $B\le1/2$, and $h_{aaaa}<4293/50=24\gamma(9/4)$ otherwise.
Since $B\ge1/3$, respectively $B\ge1/2$, this is $h_{aaaa}<24\gamma(1+B)^2$, that is, $f_{aaaa}<0$.

\emph{The point $a=1/2$.}
There, with $v=(1-K)/(1+K)$,
\[
 f_{aaa}=\frac{16q(q+1)v-16q^3v^3/r}{L}
             -\frac{54}5(1-C)-\frac{477}{25}(1-B^2).
\]
The first term increases with $v\le1/3$. Since $K^{(4)}>0$ (because $q(q-1)(q-2)(q-3)>0$),
cubic Hermite interpolation at $0,1$ gives, with $k=5247/10000<K(0)=3^{-r}$,
\[
 K\ge K_*(z)=k+(3-3k-2q/3)z^2+(2q/3-2+2k)z^3.
\]
Replacing $K$ by $K_*$ and $L$ by $L_0$, and clearing
$r(1+K_*)^3$, reduces $-f_{aaa}(1/2,z)\ge0$ to $(1-z)\sum_{j=0}^{12}c_jz^j\ge0$, where every $c_j$ is
positive (LHD-2). Lowering $K(0)$ to $k$ preserves the minorant because its coefficient is
$1-3z^2+2z^3=(1-z)^2(1+2z)\ge0$. Together with $f_{aaaa}<0$ this proves the lemma.
\end{proof}

\begin{lemma}[The whole large-head chart; exact polynomial sign certificates]\label{lem:fqe-large-head-dual}
Every spectrum on the above chart satisfies
$\Lambda_\theta<T=2507333/1500000$.
\end{lemma}
\begin{proof}
Put $a_0=127/200$, $S=1-a_0=73/200$, $d=a_0-1/2=27/200$ and $\xi=(a-a_0)/d\in[-1,11/9]$.
By Lemma~\ref{lem:fqe-large-head-curvature}, the quadratic Hermite interpolant in $a$
at the single node $1/2$ and the double node $a_0$ is a global upper bound,
since its remainder has the sign of $f_{aaa}(\eta,z)(a-1/2)(a-a_0)^2\le0$:
\[
 f(a,z)\le Q=f_0(1-\xi^2)+f_l\xi^2+dD(\xi+\xi^2),\quad
 f_0=f(a_0,z),\ f_l=f(\tfrac12,z),\ D=f_a(a_0,z).
\]
Write $f=h(a,K)+P$, where $P$ is a polynomial in $(a,B,C)$ and does not depend on $K$.
Then only the three values $h(a_0,K)$, $h(\tfrac12,K)$, $h_a(a_0,K)$ in $Q$ depend on $K$.

\emph{$Q$ decreases in $K$ on $[.52,1]$.}
Let $R=(a_0/S)^q\in[12/5,5/2]$, $c=rL$ and $\lambda=dq/(a_0S)=42849/46355$.
Differentiating,
\[
 h_K(a_0,K)=\frac{-1}{c(R+K)},\quad h_K(\tfrac12,K)=\frac{-1}{c(1+K)},\quad
 h_{aK}(a_0,K)=\frac{qR}{c\,a_0S(R+K)^2},
\]
and therefore
\[
 -c(R+K)^2(1+K)\,Q_K=(1+K)\,(R+K-\lambda R\xi)+\xi^2\{(R+K)(R-1)-\lambda R(1+K)\}.
\]
For $-1\le\xi\le11/9$ both terms are nonnegative, the first strictly positive. The first is linear in $\xi$, positive at both ends:
$\frac{11}9\lambda<1.13$ and $K>.13R$. The brace increases in $R$ and decreases in $K$, and is
$\frac{17}5\cdot\frac75-\frac{24}5\lambda>0$ at $R=12/5$, $K=1$ (LHD-3).

\emph{A cubic minorant of $K$.}
Because $K^{(4)}>0$, cubic Hermite interpolation at the double nodes $3/20,17/20$ gives
\[
 K(z)\ge\kappa(z)=.5238033+.0142244z+.4204281z^2+.0384307z^3,
 \qquad .52<\kappa<1.                                      \tag{LH1}
\]
Each coefficient of $\kappa$ is a rational lower rounding of the interpolant's coefficient (LHD-3).
Since $.52<\kappa\le K\le1$, we have $Q(K)\le Q(\kappa)$.

\emph{Rational bounds for the logarithms.} Set
\[
 \rho=2.40792594,\quad u=2.45774282,\quad H_0=1.10489073,\quad H_l=1.32818582.
\]
They approximate $R$, $1/(rL)$, $h(a_0,3/4)$, $h(\tfrac12,3/4)$ respectively with error less than $2\cdot10^{-8}$.
For $v\in[.52,1]$ define
\[
 \widehat h_0=H_0-2u\frac{v-3/4}{2\rho+v+3/4},\quad
 \widehat h_l=H_l-2u\frac{v-3/4}{v+11/4},\quad
 \widehat h_a=-\frac{2qu}{1-(27/100)^2}\Bigl(\frac{\rho-v}{\rho+v}-\frac{27}{100}\Bigr).
\]
The identity $\ln((1+y)/(1-y))=2y+e_y$ has $|e_y|\le2|y|^3/(3(1-y^2))$.
Here $h(a_0,v)-h(a_0,\tfrac34)=-\frac1{rL}\ln\frac{R+v}{R+3/4}$ and
$h(\tfrac12,v)-h(\tfrac12,\tfrac34)=-\frac1{rL}\ln\frac{1+v}{7/4}$. The logarithms equal $\ln\frac{1+y}{1-y}$ with
$y=(v-\frac34)/(2R+v+\frac34)$, respectively $y=(v-\frac34)/(v+\frac{11}4)$, and $|y|<.0383$, respectively $|y|<.0704$.
Rounding the constants and the remainder give
\[
 |h(a_0,v)-\widehat h_0|<e_0=10^{-4},\quad
 |h(\tfrac12,v)-\widehat h_l|<e_l=6\cdot10^{-4},\quad
 |h_a(a_0,v)-\widehat h_a|<e_a=10^{-6}.                         \tag{LH2}
\]
These include the rounding errors; LHD-3 checks the scalar arithmetic.
On $[-1,11/9]$ we have $|1-\xi^2|\le1$ and $|d(\xi+\xi^2)|\le11/30<1$, so with
$P_0=P(a_0)$, $P_l=P(\tfrac12)$ and $P_a=\partial_aP(a_0)$,
\[
 Q(\kappa)\le F=(\widehat h_0+P_0)(1-\xi^2)+(\widehat h_l+P_l)\xi^2+d(\widehat h_a+P_a)(\xi+\xi^2)
   +e_0+e_a+e_l\xi^2,\qquad v=\kappa(z).
\]

\emph{The completed square.}
Put $E=(v+2\rho+3/4)(v+11/4)(v+\rho)>0$. Then $E\,(T-F)=c_0+c_1\xi+c_2\xi^2$, where after
substituting $v=\kappa(z)$ and the polynomials $B,C$ the $c_i$ are polynomials in $z$.
On $[0,1]$,
\[
 c_2\ge\tfrac1{10},\qquad 4c_0c_2-c_1^2\ge\tfrac3{100},
\]
because all Bernstein coefficients of $c_2$ (degree $13$) and of $4c_0c_2-c_1^2$ (degree $26$) are at least
these values (LHD-4). Thus the quadratic in $\xi$ is positive for every real $\xi$, so $f\le Q\le F<T$.
\end{proof}
